% Figure: the stability map of a two-mirror laser cavity in the (g1, g2) plane,
% where g_i = 1 - L/R_i. A resonator is stable when the round-trip ray-transfer
% matrix has |(A+D)/2| <= 1, which reduces to 0 <= g1 g2 <= 1. The stable region
% lies between the two branches of the hyperbola g1 g2 = 1 and the axes. Named
% configurations sit at marked points: plane-parallel (1,1), confocal (0,0),
% concentric (-1,-1), and the symmetric hemispherical case.
\begin{figure}[htbp]
\centering
\begin{tikzpicture}
\begin{axis}[
  width=10.5cm, height=9cm,
  xlabel={$g_1 = 1 - L/R_1$},
  ylabel={$g_2 = 1 - L/R_2$},
  xmin=-1.6, xmax=1.6, ymin=-1.6, ymax=1.6,
  axis lines=middle,
  xtick={-1,1}, ytick={-1,1},
  tick label style={font=\scriptsize},
  label style={font=\small},
  samples=200,
]
% stable region boundary: g1 g2 = 1, two hyperbola branches
\addplot[engblue, very thick, domain=0.625:1.6] {1/x};
\addplot[engblue, very thick, domain=-1.6:-0.625] {1/x};
% shade a representative stable cell in the first quadrant under g1 g2 = 1
\addplot[engblue!12, domain=0.001:1, samples=120] {min(1/x,1.6)} \closedcycle;
\addplot[engblue!12, domain=-1:-0.001, samples=120] {max(1/x,-1.6)} \closedcycle;
% named configurations
\addplot[only marks, mark=*, engred, mark size=1.6pt]
  coordinates {(1,1) (0,0) (-1,-1) (0,1) (1,0)};
\node[engred, font=\scriptsize, anchor=south west] at (axis cs:1.0,1.02)
  {plane-parallel};
\node[engred, font=\scriptsize, anchor=south west] at (axis cs:0.05,0.05)
  {confocal};
\node[engred, font=\scriptsize, anchor=north east] at (axis cs:-1.0,-1.02)
  {concentric};
\node[engred, font=\scriptsize, anchor=south] at (axis cs:0,1.05)
  {hemispherical};
\node[engblue, font=\scriptsize, anchor=west] at (axis cs:0.55,0.78)
  {$g_1 g_2 = 1$};
\node[enggray, font=\scriptsize, anchor=center] at (axis cs:0.45,0.42)
  {stable};
\node[enggray, font=\scriptsize, anchor=center] at (axis cs:-0.55,-0.55)
  {stable};
\end{axis}
\end{tikzpicture}
\caption[Laser-cavity stability map]{Stability of a two-mirror optical
resonator in the plane of the $g$-parameters $g_i = 1 - L/R_i$, with $L$ the
mirror spacing and $R_i$ the mirror radii. The round-trip ray-transfer matrix
returns a ray to the neighbourhood of its starting line only when its trace
satisfies $|(A+D)/2|\le 1$, which reduces to the resonator condition
$0 \le g_1 g_2 \le 1$. The shaded cells between the axes and the hyperbola
$g_1 g_2 = 1$ are the stable configurations; outside them a ray walks off the
mirrors after a finite number of bounces and the cavity does not support a
mode. Named designs sit at the corners and edges: plane-parallel at $(1,1)$,
confocal at the origin, concentric at $(-1,-1)$, each on the boundary where
the stability margin vanishes, which is why working lasers are designed a
little inside an edge rather than on it.}
\label{fig:vol04ch11:abcd-cavity-stability}
\end{figure}
